\section{Magnetic Field of a Single Infinite Sheet (Derivation)}

Continuing from the infinite-sheet calculation. The field from a single infinite current sheet with surface current density $K$ (in the $x$-direction):

\[
    B = \int_{-w/2}^{w/2} \frac{\mu_0 I}{2\pi r_w} \, dx \cdot \cos\theta
\]

\[
    B = \int_{-w/2}^{w/2} \frac{\mu_0 I \, y}{2\pi w (x^2 + y^2)} \, dx
\]

Using the standard integral $\displaystyle\int_{-a}^{a} \frac{1}{x^2 + y^2} \, dy = \frac{\pi}{y}$, so:

\[
    \Rightarrow \quad B = \frac{\mu_0 I}{2w} \, \hat{x}
\]

\section{Magnetic Field of Two Parallel Planes}

\begin{figure}[H]
\begin{center}
\begin{tikzpicture}[scale=0.9]
  % Two sheets
  \fill[orange!15] (-3.5,0.9) rectangle (3.5,1.15);
  \fill[orange!15] (-3.5,-1.15) rectangle (3.5,-0.9);
  \draw[thick, orange!70!black] (-3.5,1.0) -- (3.5,1.0) node[right]{$K\hat{x}$};
  \draw[thick, orange!70!black] (-3.5,-1.0) -- (3.5,-1.0) node[right]{$-K\hat{x}$};
  % Between sheets: B fields add (pointing in +y direction, shown as arrows)
  \foreach \x in {-2.5,-1.5,-0.5,0.5,1.5,2.5} {
    \draw[-{Stealth[length=4pt]}, thick, blue!70!black]
      (\x, -0.85) -- (\x, 0.85);
  }
  \node[blue!70!black] at (0,0) {$\vect{B} = \mu_0 K\hat{y}$};
  % Outside: fields cancel
  \node[gray, font=\small] at (0,1.55) {$\vect{B} = 0$ (outside)};
  \node[gray, font=\small] at (0,-1.55) {$\vect{B} = 0$ (outside)};
\end{tikzpicture}
\end{center}
\caption{Two infinite parallel current sheets with opposite surface current densities $+K\hat{x}$ and $-K\hat{x}$. Between the sheets the magnetic fields add ($B = \mu_0 K$); outside they cancel ($B = 0$). This is the magnetic analogue of a parallel-plate capacitor.}
\end{figure}

The enclosed current is $I = \int \vect{J} \cdot d\vect{S}$, and Amp\`{e}re's law in differential form is:
\[
    \curl \vect{B} = \mu_0 \vect{J}
\]

\subsection{Amp\`{e}re's Law for an Infinite Sheet}

Draw an Amp\`{e}rian loop of length $L$ straddling the sheet:
\[
    2BL = \mu_0 K L \quad \Rightarrow \quad B = \frac{\mu_0 K}{2}
\]

\subsection{Solenoid}

Using an Amp\`{e}rian loop of length $L$ inside and outside the solenoid (outside $B=0$, inside $\oint \vect{B}\cdot d\vect{\ell} = \mu_0 I_{\text{enc}}$):
\[
    \mu_0 I n L = BL \quad \Rightarrow \quad B = \mu_0 n I
\]
where $n$ is the number of turns per unit length.

\begin{figure}[H]
\begin{center}
\begin{tikzpicture}[scale=0.9]
  % Solenoid body (series of ellipses for coil)
  \foreach \i in {0,1,2,3,4,5,6,7} {
    \draw[thick, blue!60!black] ({0.5*\i},0) ellipse (0.18 and 0.6);
  }
  % Axis
  \draw[dashed, gray] (-0.5,0) -- (4.5,0);
  % Uniform B field inside
  \foreach \y in {-0.3,0,0.3} {
    \draw[-{Stealth[length=4pt]}, thick, red!70!black]
      (0.2,\y) -- (3.3,\y);
  }
  \node[red!70!black] at (1.75,0.5) {$\vect{B} = \mu_0 nI$};
  % Outside B = 0
  \node[gray, font=\small] at (4.2,0.8) {$B=0$};
  % Amperian rectangle
  \draw[dashed, green!60!black, thick]
    (0.6,-0.9) rectangle (3.0,0.9);
  % Amperian loop label
  \node[green!60!black, font=\small] at (1.8,-1.2) {Amp\`{e}rian loop ($L$)};
  % Current I annotation
  \node[font=\small] at (-0.3,-0.8) {$I$};
\end{tikzpicture}
\end{center}
\caption{Solenoid with $n$ turns per unit length. The Amp\`{e}rian rectangle (dashed) of length $L$ has one side inside (field $B$) and one outside ($B=0$). Amp\`{e}re's law gives $BL = \mu_0 nIL$, so $B = \mu_0 nI$.}
\end{figure}

\section{Hall Effect}

\begin{figure}[H]
\begin{center}
\begin{tikzpicture}[scale=0.9]
  % Conductor body
  \fill[yellow!15] (0,-0.8) rectangle (4.5,0.8);
  \draw[thick] (0,-0.8) rectangle (4.5,0.8);
  % Current to the right
  \foreach \y in {-0.4,0,0.4} {
    \draw[-{Stealth[length=4pt]}, thick, red!70!black]
      (0.5,\y) -- (1.8,\y);
  }
  \node[red!70!black] at (1.1,0.55) {$I$};
  % Magnetic field B downward (into page, shown as crosses)
  \foreach \x in {0.5,1.5,2.5,3.5,4.0} {
    \draw[gray!60] (\x,1.4) circle (0.12) node[font=\tiny]{$\times$};
  }
  \node[gray, font=\small] at (2.0,1.8) {$\vect{B}$ (into page)};
  % Charge accumulation on bottom
  \foreach \x in {1.0,2.0,3.0,4.0} {
    \node[blue!70!black, font=\small] at (\x,-0.65) {$-$};
    \node[red!70!black, font=\small] at (\x,0.65) {$+$};
  }
  % Hall voltage arrow
  \draw[<->, thick, orange!80!black]
    (4.8,-0.8) -- (4.8,0.8) node[midway, right]{$V_H$};
  % Hall E field arrow
  \draw[-{Stealth[length=5pt]}, thick, blue!60!black]
    (2.3,-0.5) -- (2.3,0.5) node[right, font=\small]{$E_H$};
\end{tikzpicture}
\end{center}
\caption{Hall effect: current $I$ flows right in a conductor; field $\vect{B}$ is directed into the page. Charge carriers deflect downward (for electrons), building up $E_H$ until $eE_H = evB$. The Hall voltage $V_H = BIh/(neh\ell) = BI/(ne\ell)$ measures charge carrier density $n$.}
\end{figure}

In the Hall effect, a transverse magnetic field deflects current carriers to one side of the conductor, building up a transverse electric field that balances the magnetic force. In steady state the electric force exactly cancels the magnetic force:
\[
    eE_H = ev B \quad \Rightarrow \quad E_H = vB.
\]
The current density $\vect{J} = -ne\vect{v}$ gives the drift speed $v = J/(ne) = I/(neh\ell)$, and the Hall voltage $V_H = E_H \ell$ becomes:
\[
    V_H = \left(\frac{B}{neh}\right) I.
\]
The sign of the Hall voltage reveals the sign of the charge carriers, and its magnitude measures the carrier density $n$.

\section{Maxwell's Equations So Far}

\begin{formula}[Maxwell's Equations (magnetostatics)]
\begin{align*}
    \div \vect{E} &= \frac{\rho}{\varepsilon_0} & \div \vect{B} &= 0 \\
    \curl \vect{E} &= 0 & \curl \vect{B} &= \mu_0 \vect{J}
\end{align*}
\end{formula}

\section{Magnetic Vector Potential}

Let $\vect{B} = \curl \vect{A}$ (the \emph{magnetic vector potential}). Then $\div\vect{B} = 0$ is automatically satisfied because $\div(\curl \vect{A}) = 0$. The vector potential $\vect{A}$ is not uniquely defined: adding any gradient $\grad f$ to $\vect{A}$ leaves $\vect{B}$ unchanged (gauge freedom). The Coulomb gauge $\div\vect{A}=0$ is the standard choice in magnetostatics; it leads to $\laplacian\vect{A} = -\mu_0\vect{J}$, a vector Poisson equation analogous to the scalar Poisson equation $\laplacian\phi = -\rho/\epsilon_0$.

In component form:
\[
    \vect{B} = \begin{vmatrix} \hat{x} & \hat{y} & \hat{z} \\ \partial_x & \partial_y & \partial_z \\ A_x & A_y & A_z \end{vmatrix}
    \quad \Rightarrow \quad B_x = \frac{\partial A_z}{\partial y} - \frac{\partial A_y}{\partial z}, \quad \text{etc.}
\]

Example: for a uniform field $\vect{B} = B_0 \hat{z}$,
\[
    \vect{A} = B_0 x \hat{y} \quad \text{or} \quad \vect{A} = -B_0 y \hat{x}
\]

\subsection{Is \texorpdfstring{$\vect{A}$}{A} Unique? (Gauge Freedom)}

We also know $\curl(\grad f) = 0$ for any scalar $f$. So if $\vect{A}$ and $\vect{A}'$ both give the same $\vect{B}$:
\[
    \curl(\vect{A} - \vect{A}') = 0 \quad \Rightarrow \quad \vect{A} - \vect{A}' = \grad\psi
\]
for some scalar field $\psi$.

We can fix the gauge by imposing $\div \vect{A} = 0$ (Coulomb gauge):
\[
    \vect{A} - \vect{A}' = \grad\psi \quad \Rightarrow \quad \div\vect{A} - \div\vect{A}' = \nabla^2 \psi
\]

Suppose we have $\curl\vect{A} = \vect{B}$ and $\div\vect{A} = f(\vect{r})$. We know:
\[
    \div\vect{E} = \frac{\rho(x,y,z)}{\varepsilon_0}
\]
We can find a vector field such that $\div\vect{A}' = f(\vect{r})$, then subtract it.

\begin{fact}
Be careful about integrating electric/magnetic fields from scratch: you could get a wrong answer automatically by not considering volume currents infinitely close to the point in question (see exercise 1.69, pg.\ 305). Solution: remove a thin strip.
\end{fact}

We already know for electrostatics:
\[
    \vect{E} = -\grad\phi, \qquad \nabla^2\phi = -\frac{\rho}{\varepsilon_0}
\]

\subsection{Uniform Magnetic Field and the Vector Potential Poisson Equation}

In the Coulomb gauge ($\div\vect{A}=0$), Amp\`{e}re's law becomes:
\[
    \curl\vect{B} = \mu_0\vect{J}
\]
\[
    \curl(\curl\vect{A}) = \mu_0\vect{J}
\]
\[
    = \grad(\div\vect{A}) - \nabla^2\vect{A} = -\nabla^2\vect{A}
\]
since $\div\vect{A} = 0$. This gives three decoupled Poisson equations:

\begin{formula}[Poisson Equations for the Vector Potential]
\[
    \nabla^2 A_x = -\mu_0 J_x, \qquad \nabla^2 A_y = -\mu_0 J_y, \qquad \nabla^2 A_z = -\mu_0 J_z
\]
\end{formula}

These are exactly analogous to $\nabla^2\phi = -\rho/\varepsilon_0$.

\section{Solution for \texorpdfstring{$\vect{A}$}{A}: Analogy with Electrostatics}

Recall that:
\[
    \phi(\vect{r}_1) = \iiint \frac{\rho_2}{4\pi\varepsilon_0 r_{12}} \, dV_2
\]

So for $A_x$:
\[
    A_x = \frac{\mu_0}{4\pi} \int \frac{J_x}{r_{12}} \, dV
\]

In general:
\begin{formula}[Magnetic Vector Potential from Current]
\[
    \vect{A}(\vect{r}_1) = \frac{\mu_0}{4\pi} \int \frac{\vect{J}(\vect{r}_2)}{r_{12}} \, dV
\]
\end{formula}

\subsection{Example: Long Wire}

By analogy with the cylindrical capacitor ($E_r = \frac{\lambda}{2\pi\varepsilon_0 r}$, $\phi = \frac{\lambda}{2\pi\varepsilon_0}\ln(a/r)= \frac{\rho a^2}{2\varepsilon_0}\ln(a/r)$):

\[
    A_x = \frac{J_x \mu_0 a^2}{2} \ln\!\left(\frac{a}{r}\right) = \frac{\mu_0 I}{2\pi} \ln\!\left(\frac{a}{r}\right)
\]

So $\vect{A} = (A_x, 0, 0)$ and $\vect{B} = \curl\vect{A}$.

\[
    \vect{B} = -\frac{\partial A_x}{\partial z}\,\hat{z} + \cdots
\]
\[
    \frac{\partial A_x}{\partial z} = \left(\frac{-\mu_0 I z}{2\pi r^2}\right) = B_y, \quad \text{etc.}
\]

\section{Biot--Savart Law from the Vector Potential}

Biot-Savart gives the magnetic field produced by a steady current element $Id\vect{\ell}$ at displacement $\vect{r}$. Note the cross product: the field circles around the current element, unlike the radial Coulomb field.

\[
    d\vect{B} = \curl(d\vect{A}), \quad \text{where} \quad d\vect{A} = \frac{\mu_0 I}{4\pi r}\,d\vect{\ell}
\]
\[
    d\vect{B} = \frac{\mu_0 I}{4\pi}\,\curl\!\left(\frac{d\vect{\ell}}{r}\right)
\]

Using the product rule $\curl\!\left(\frac{d\vect{\ell}}{r}\right) = \frac{1}{r}\,\curl(d\vect{\ell}) + \grad\!\left(\frac{1}{r}\right)\times d\vect{\ell}$, and noting that $d\vect{\ell}$ is constant so $\curl(d\vect{\ell})=0$, and $\grad(1/r) = -\hat{r}/r^2$:

\[
    d\vect{B} = \frac{\mu_0 I}{4\pi}\left(d\vect{\ell} \times \grad\!\left(\frac{1}{r}\right)\right), \qquad \grad\!\left(\frac{1}{r}\right) = -\frac{\hat{r}}{r^2}
\]

\begin{formula}[Biot--Savart Law]
\[
    d\vect{B} = \frac{\mu_0 I}{4\pi r^2}\, d\vect{\ell} \times \hat{r}
\]
\end{formula}

\subsection{Field on Axis of a Circular Loop}

\begin{figure}[H]
\begin{center}
\begin{tikzpicture}[scale=1.0]
  % Circular loop (ellipse to suggest 3D)
  \draw[thick, blue!70!black] (0,0) ellipse (1.8 and 0.5);
  % Axis (z)
  \draw[->] (0,-1.8) -- (0,2.8) node[right]{$z$};
  % Field point P on z-axis
  \filldraw[black] (0,2.0) circle (3pt) node[right]{$P$};
  % z label
  \draw[<->, gray] (0.25,0) -- (0.25,2.0) node[midway,right, font=\small]{$z$};
  % R label
  \draw[<->, gray] (0,0.08) -- (1.8,0.08) node[midway,above, font=\small]{$R$};
  % Line element dl on loop
  \draw[-{Stealth[length=5pt]}, thick, red!70!black]
    (1.6,0.3) -- (1.0,0.48)
    node[above right, font=\small]{$d\vect{\ell}$};
  % r vector from dl to P
  \draw[dashed, orange!80!black, -{Stealth[length=4pt]}]
    (1.6,0.3) -- (0,2.0) node[midway, right, font=\small]{$r$};
  % dB direction (along z at axis by symmetry)
  \draw[-{Stealth[length=5pt]}, thick, green!60!black]
    (0,2.0) -- (0,2.7) node[right, font=\small]{$d\vect{B}$};
  % Loop label
  \node[blue!60!black, font=\small] at (-2.2,0) {current $I$};
\end{tikzpicture}
\end{center}
\caption{Circular current loop of radius $R$ in the $xy$-plane. The field point $P$ is on the $z$-axis at height $z$. By symmetry, transverse components of $d\vect{B}$ cancel, and the net field is along $\hat{z}$.}
\end{figure}

\[
    B = \frac{\mu_0 I}{4\pi r^2} \int \frac{d\vect{\ell}\times\hat{r}}{r}
\]
\[
    \vect{B} = \frac{\mu_0 I}{4\pi} \int \frac{R\sin\theta}{R^2 + z^2} \, d\phi
\]

\begin{formula}[Field on Axis of a Circular Loop]
\[
    \vect{B} = \frac{\mu_0 I}{2} \frac{R^2}{(R^2 + z^2)^{3/2}}\,\hat{z}
\]
\end{formula}

\subsection{Field on Axis of a Finite Solenoid}

Each ring of the solenoid contributes:
\[
    d\vect{B} = \frac{\mu_0 R^2}{2(R^2+z^2)^{3/2}} n I \, dz
\]
Integrating from $-L/2$ to $L/2$:
\[
    \vect{B} = \frac{\mu_0 n I}{2} \int_{-L/2}^{L/2} \frac{R^2}{(R^2+z^2)^{3/2}} \, dz
\]
\[
    \vect{B} = \mu_0 n I \frac{L}{\sqrt{L^2 + 4R^2}}\,\hat{z}
\]


\section{Lorentz Force}

\begin{formula}[Lorentz Force]
\[
    \vect{F} = q\left(\vect{E} + \vect{v}\times\vect{B}\right)
\]
\end{formula}

For a long straight wire, the field is:
\[
    \vect{B} = \frac{\mu_0 I}{2\pi r}\left(\hat{\imath}\times\hat{r}\right) = \frac{\mu_0 I}{2\pi r}\,\hat{\phi}
\]

\subsection{Example: Wien's Filter (Velocity Selector)}

Force balance:
\[
    \vect{F} = q(\vect{E} + \vect{v}\times\vect{B}) = q(E + vB) = 0
\]
\[
    \Rightarrow \quad E = -vB \quad \Rightarrow \quad v = \frac{E}{B}
\]

The Wien filter balances electric and magnetic forces so that only particles with $v = E/B$ pass straight through. It selects a single speed regardless of the charge or mass of the particle.


\section{Force on a Current-Carrying Wire}

The force on a current element $dq\,\vect{v} = dq\,\frac{d\vect{s}}{dt} = I\,d\vect{s}$ in a field $\vect{B}$:
\[
    \vect{F} = \int I\,d\vect{s}\times\vect{B}
\]

\subsection{Example: Current Loop Levitated by Magnetic Force}

Equilibrium condition $\vect{F}_g + \vect{F}_B = 0$ with $\vect{F}_g = m\vect{g} = -mg\,\hat{z}$:

\[
    \vect{F}_B = \int I\,d\vect{s}\times\vect{B} = \int_0^\ell I\,d\vect{s}\times B\,\hat{x} = \int_0^\ell I\,dy\,B\,\hat{z} = IB\ell\,\hat{z}
\]

Thus $IB\ell = mg$, giving:
\[
    \vect{I} = \frac{mg}{B\ell}\,\hat{y}
\]

\section{Example: Field of a Plane Current Sheet}

A single infinite wire gives $\vect{B} = \frac{\mu_0 J}{2\pi r}\,\hat{\phi}$.

In the continuum limit (many wires $\to$ a sheet), the surface current density is:
\[
    K = \text{line current density} = \frac{I}{\ell}
\]
The current in a strip of width $dx$ is $dI = K\,dx$.

\[
    dB_x = \frac{\mu_0 I}{2\pi r} \cdot \frac{K\,dx}{r} \cdot \cos\theta, \qquad \cos\theta = \frac{y}{r}
\]
\[
    dB_x = \frac{\mu_0 K}{2\pi} \cdot \frac{y\,dx}{x^2+y^2}
\]
\[
    B_x = \int_{-\infty}^{\infty} \frac{\mu_0 K y \, dx}{2\pi(x^2+y^2)} = \frac{\mu_0 K y}{2|y|}\,\hat{x}
\]

Thus:
\[
    \vect{B} = \begin{cases} \dfrac{\mu_0 K}{2}\,\hat{x} & y \geq 0 \\[6pt] -\dfrac{\mu_0 K}{2}\,\hat{x} & y \leq 0 \end{cases}
\]

The infinite sheet is independent of $y$ (distance from the sheet). When we cross the sheet of current:
\[
    |\Delta\vect{B}| = \mu_0 K
\]
This is similar to crossing a sheet of charge, where $|\Delta\vect{E}| = \sigma/\varepsilon_0$, but the $B$ field is parallel (not perpendicular) to the sheet.

\subsection{Amp\`{e}re's Law Verification}

Using a rectangular Amp\`{e}rian loop of width $\ell$ straddling the sheet:
\[
    \oint \vect{B}\cdot d\vect{\ell} = I_{\text{enc}}\,\mu_0
\]
\[
    2B\ell = K\ell \quad \Rightarrow \quad |\vect{B}| = \frac{\mu_0 K}{2}
\]

\section{How the Fields Transform}

\subsection{Field Transformation Under a Lorentz Boost}

\subsection{Summary of Field Transformations}

\begin{formula}[Electromagnetic Field Transformations]
\begin{align*}
    E'_\| &= E_\| & E'_\perp &= \gamma\left(E_\perp + \vect{v}\times\vect{B}_\perp\right) \\[4pt]
    B'_\| &= B_\| & B'_\perp &= \gamma\left(\vect{B}_\perp - \frac{\vect{v}}{c^2}\times\vect{E}_\perp\right)
\end{align*}
\end{formula}

where $\|$ and $\perp$ refer to directions parallel and perpendicular to the boost velocity $\vect{v}$.

\chapter{Electromagnetic Induction}

\section{Conducting Rod in a Magnetic Field}

The charges inside the moving rod experience a force $\vect{F} = q\vect{v}\times\vect{B}$ (not a force from $\vect{E}$). This creates a charge separation (positive charges accumulate at one end).

\subsection{Electromotive Force (EMF)}

\begin{definition}[Electromotive Force]
\[
    \mathcal{E} \equiv \frac{1}{q}\oint \vect{f}\cdot d\vect{s}
\]
The EMF is the electromotive force (not actually a force --- it has units of electric potential). It is any influence that causes charge to circulate around a closed path. EMF must be non-electrostatic, since $\curl\vect{E} = 0$ for electrostatic fields and a conservative field cannot drive a current around a closed loop.
\end{definition}

Electromagnetic induction arises when the magnetic flux through a circuit changes, whether by moving the circuit, changing the current in a nearby circuit, or changing $\vect{B}$ in time. The induced EMF is $\mathcal{E} = -d\Phi_B/dt$ (Faraday's law).

\begin{fact}[Useful Unit Conversion]
\[
    \frac{1\,\text{m}^2 \cdot \text{Tesla}}{\text{second}} = 1\,\text{Volt}
\]
\end{fact}

